\begin{abstract}
Examination-level multimodal multi-label recognition must combine an ordered
set of images with a paired report, while different findings may depend on
different images and textual evidence. Existing methods often encode sampled
images with fixed index distances and fuse image--report features into a shared
representation, which can distort contextual relationships and mix evidence
across labels. We present Acquisition-aware Label-centric Multimodal Multiple
Instance Learning (ALM-MIL), built on Acquisition-aware Contextual Position
Encoding (ACPE) and Label-Centric Coexistence Fusion (LCCF). ACPE uses
acquisition order as a structural anchor and learns content-adaptive contextual
intervals from neighboring visual transitions, injecting the resulting
coordinates through complementary absolute and relative routes. LCCF organizes
evidence around individual labels by selecting label-specific visual instances,
modeling label coexistence with hypergraph reasoning, retrieving
label-conditioned textual evidence, and adaptively fusing the retrieved
information. Experiments on four clinical cohorts and four public benchmarks
demonstrate the superiority of ALM-MIL over existing models. Further
experiments on structured image missingness and
label-centric evidence organization validate ACPE's ability to model
acquisition-aware contextual relations and LCCF's ability to organize
label-specific multimodal evidence.

\cntranslation{%
检查级多模态多标签识别需要联合处理一组有序图像和一份配对报告，而不同异常可能依赖不同的图像和文本证据。现有方法通常使用固定的采集索引间距对采样图像进行编码，并将图像与报告特征融合为共享表征，容易造成上下文关系失真，并混合不同标签之间的证据。为此，我们提出采集感知的标签中心多模态多实例学习框架（Acquisition-aware Label-centric Multimodal Multiple Instance Learning，ALM-MIL），其核心由采集感知上下文位置编码（Acquisition-aware Contextual Position Encoding，ACPE）和标签中心共存融合（Label-Centric Coexistence Fusion，LCCF）组成。ACPE 以采集顺序作为结构参照，根据相邻图像的视觉变化学习内容自适应的上下文间隔，并通过互补的绝对路径和相对路径注入相应坐标。LCCF 围绕各个标签组织证据，通过选择标签特异性的视觉实例、利用超图推理建模标签共存关系、检索标签条件文本证据，并自适应融合检索结果，实现标签级多模态建模。我们在四个临床队列和四个公开基准上的实验验证了 ALM-MIL 相比现有模型的优越性。针对结构化图像缺失和标签中心证据组织的进一步实验，验证了 ACPE 建模采集感知上下文关系的能力，以及 LCCF 组织标签特异性多模态证据的能力。
}
\end{abstract}
